\section{Design}
\label{sec:design}

ActPlane is a programmable OS-level control plane for agent harnesses. It
compiles behavioral policies into labeled information-flow rules at OS
kernel hooks, observing and enforcing agent behavior and returning semantic
feedback when a rule is matched. Figure~\ref{fig:architecture} summarizes
the pipeline.

\begin{figure}[t]
\begin{verbatim}
actplane.yaml
  policy: |
        |
        v
Rust compiler ------> struct taint_config
        |                    |
        v                    v
 actplane run -----> eBPF loader + BPF maps
        |                    |
        v                    v
  target command <--- kernel hooks
        |
        v
.actplane/last-violation.txt
\end{verbatim}
\caption{ActPlane pipeline. The Rust compiler lowers the policy DSL into a
fixed-size binary configuration. The eBPF loader writes it into read-only
data, attaches to kernel hooks, and reports rule matches as NDJSON. The runner
maps rule identifiers back to human-readable feedback.}
\Description{ASCII diagram of policy compilation, eBPF loading, target
command execution, and feedback output.}
\label{fig:architecture}
\end{figure}

\subsection{Policy Model}

ActPlane policies are object-level information-flow rules over three node
types: processes (identified by pid and command name), files (keyed by
path hash), and endpoints (keyed by IPv4 address and mask). The model
operates at object granularity rather than byte
granularity~\cite{kemerlis2012libdft}, trading precision for low-overhead
whole-system coverage.

\subsubsection{Labels and Propagation}

Each node carries a 64-bit label bitmask. A \emph{source} declaration
introduces labels: an exec source labels processes that execute a matching
binary, a file source labels files matching a path pattern, and an endpoint
source labels network endpoints matching an address range. Labels propagate
through fixed transfer functions:

Let $L(x)$ denote the label bitmask of node $x$. Labels propagate through
fixed transfer functions:

\begin{itemize}
\item \textbf{fork}$(p \to c)$: $L(c) := L(c) \cup L(p)$.
\item \textbf{exec}$(p, f)$: $L(p) := L(p) \cup L(f) \cup S(f)$, where
$S(f)$ is the set of source labels matching $f$.
\item \textbf{read}$(p, f)$: $L(p) := L(p) \cup L(f)$.
\item \textbf{write}$(p, f)$: $L(f) := L(f) \cup L(p)$.
\item \textbf{connect}$(p, e)$: $L(e) := L(e) \cup L(p)$.
\end{itemize}

All transfer functions are monotonic: labels accumulate and are never
removed by propagation. The only mechanism that removes a label is an
explicit \emph{declassify} directive in the DSL, which clears specified
labels when the agent executes a designated gate binary (e.g., a redaction
tool). This monotonicity bounds over-tainting: a node's label set grows
at most to the union of all source labels reachable through observed flow
edges.

These transfer functions maintain a provenance invariant: if a process $p$
carries label $\ell$, then information from some source tagged with $\ell$
has reached $p$ through observed OS-level flow edges.

\subsubsection{Rules and Effects}

A rule contains one or more clauses. Each clause starts with an action verb
(\texttt{notify}, \texttt{block}, or \texttt{kill}) followed by an
operation (\texttt{exec}, \texttt{read}, \texttt{write}, \texttt{unlink},
\texttt{connect}, or \texttt{open}), a target pattern, optional positional
arguments, a Boolean expression over required and forbidden labels, and an
optional condition (target allow-list, lineage gate, or temporal after-gate).

\begin{figure}[t]
\begin{verbatim}
source AGENT = exec "**/claude"
source SECRET = file "**/.env"

rule no-secret-egress:
  block connect endpoint "*" if SECRET
  because "secret-derived data must stay local; run the redactor first"

rule no-agent-git:
  kill exec "git" if AGENT
  because "agent must not invoke git directly"
\end{verbatim}
\caption{Two ActPlane policy rules. The first uses BPF-LSM pre-operation
denial; the second uses harness-level termination. Exec target patterns
without \texttt{/} are implicitly basename-matched.}
\Description{Code listing of two ActPlane rules.}
\label{fig:policy-example}
\end{figure}

Each clause declares its own effect that determines how the kernel
responds to a match.

\textbf{Notify}: the kernel emits a match event and allows the operation
to proceed.

\textbf{Block}: a BPF-LSM hook returns \texttt{-EPERM} before the operation
commits. Block is available only when the \texttt{bpf} LSM is active in the
running kernel; it is not silently downgraded.

\textbf{Kill}: the kernel sends \texttt{SIGKILL} to the matching task. From
a BPF-LSM hook, kill also returns \texttt{-EPERM}; from a tracepoint, kill
is post-operation termination.

When multiple rules match the same event, the strongest effect applies in
the order kill $>$ block $>$ notify.

\subsubsection{Policy Language Scope}

The DSL is not a general-purpose mandatory access control policy. Its
primitives map to recurring agent workflow patterns: ``after tests,''
``through a review tool,'' ``data derived from untrusted input,'' and
``read-only sub-agent.'' The grammar is small so that policy authors can
translate natural-language instructions into source/sink/gate shapes without
security-policy expertise.

\subsection{System Architecture}

\subsubsection{Compiler and Kernel ABI}

The Rust compiler parses the DSL and lowers it into a C-compatible
\texttt{struct taint\_config}. Lowering allocates label bits, expands Boolean
expressions into required and forbidden bitmasks, lowers glob patterns into
exact, prefix, suffix, or wildcard match kinds, and compiles IPv4 patterns
into network/mask pairs. Human-readable rule names and reason strings
(including remediation guidance) remain in Rust metadata; the kernel receives
only fixed-size rule tables and emits integer rule identifiers.

This split keeps the eBPF program verifier-friendly. The kernel program
iterates over fixed-size arrays in read-only data using bounded loops, copies
each rule entry into a stack-local variable before matching, and uses explicit
index bounds checks for argument scanning and pattern comparison.

\subsubsection{Kernel State}

The eBPF backend maintains five BPF hash maps:

\begin{itemize}
\item \texttt{ts\_proc}: pid $\to$ label bitmask and lineage-gate state.
\item \texttt{ts\_root}, \texttt{ts\_sess}: root-process and session-level
temporal (after) gate state.
\item \texttt{ts\_file}: path hash $\to$ label bitmask.
\item \texttt{ts\_endp}: IPv4 address $\to$ label bitmask.
\end{itemize}

BPF map entries are keyed by tid (one entry per Linux thread) for
\texttt{ts\_proc} and by path hash or IPv4 address for object maps.
Updates use atomic BPF helper operations; concurrent reads from different
CPUs observe a consistent label bitmask per entry.

Tracepoint programs handle process lifecycle events (fork, exec, exit) and
post-operation file and connect events. BPF-LSM programs attach to
\texttt{bprm\_check\_security}, \texttt{file\_open}, and
\texttt{socket\_connect} for pre-operation enforcement. When BPF-LSM is not
active, the loader uses tracepoint mode and disables block-effect rules.
In tracepoint mode, hooks fire after the operation commits; a
\texttt{connect} tracepoint observes an already-established connection.
This is why \texttt{block} requires BPF-LSM: only LSM hooks
provide a pre-operation verdict. Tracepoint mode supports \texttt{notify}
(post-operation observation) and \texttt{kill} (post-operation task
termination).

The current implementation supports up to 32 sources, 32 rules, 16
transforms, and 16 gates per policy. Rule evaluation uses
\texttt{bpf\_loop} for bounded iteration, which the eBPF verifier accepts
for policies exceeding 100 rules.

\subsubsection{Unified Event Handling}

Each hook normalizes kernel arguments into a common event structure:
subject pid, object kind, access kind, backend mode, target string,
optional argument vector, and optional IPv4 address. A shared handler
resolves the subject's label bitmask, applies source labels for
read/connect events, evaluates all supported rules, emits at most one
\texttt{TAINT\_VIOLATION} event for the highest-priority matching rule, and
then applies propagation updates.

\subsection{Corrective Feedback}

The kernel match event contains the rule identifier, effect, target
string, pid, and flags indicating whether the operation was blocked or the
task was killed. The Rust runner maps the rule identifier to the
corresponding rule name and reason string, then appends a
structured payload to \texttt{.actplane/last-violation.txt}.

Agent integration uses a post-tool hook. The \texttt{actplane feedback-hook}
subcommand reads new bytes from the feedback file and returns them as
\texttt{additionalContext} to compatible agent
runtimes~\cite{agentsight}. The hook is a read-only relay: it does not
re-evaluate policy in userspace. The kernel is the sole authority for rule
matches; userspace formats the result so the agent can select an alternative
execution path.


\section{Implementation}
\label{sec:implementation}

ActPlane is implemented as a Rust driver and an eBPF program pair. The
system builds a single \texttt{actplane} binary that embeds the compiled
eBPF loader, extracts it at runtime, compiles the project's policy DSL,
and runs a target command under the observation and enforcement harness.

\subsection{Runner and Policy Discovery}

The command \texttt{actplane run -- <cmd>} discovers
\texttt{actplane.yaml} by searching upward from the working directory.
The policy YAML contains a \texttt{policy:~|} DSL block. Before the
target process begins executing, the loader seeds the stopped target's pid
with the \texttt{AGENT} label so that all descendants inherit the agent
identity. The runner prepares the feedback file
(\texttt{.actplane/last-violation.txt}) and a hook state file, exports
their paths to the child environment, waits for the loader to report
readiness, then resumes the target. This sequencing ensures that the
first target operation is observed.

\subsection{BPF Backends}

The eBPF program attaches to process lifecycle tracepoints (fork, exec,
exit), syscall tracepoints for file and connect events, and BPF-LSM hooks
for pre-operation enforcement. At load time, the loader checks whether the
running kernel has the \texttt{bpf} LSM active in
\texttt{/sys/kernel/security/lsm}. If not, it falls back to tracepoint
mode, in which the supported effects are \texttt{notify} and \texttt{kill};
\texttt{block} rules are disabled because there is no pre-operation verdict
path.

\subsection{Limitations}

\textbf{Path identity.} File labels are keyed by FNV-1a path hash rather
than inode/device pairs. Hard links and mount binds can cause the same
file to carry different labels under different paths.

\textbf{Endpoint identity.} Connect matching uses numeric IPv4 addresses.
Hostname, DNS, SNI, and HTTP-layer matching require userspace resolution or
additional instrumentation.

\textbf{Argument predicates.} Positional argument predicates inspect
arguments after exec via the tracepoint argument buffer. Pre-operation
blocking of argument predicates requires an argument cache before
\texttt{bprm\_check\_security}, which is not yet implemented.

\textbf{Label precision.} Labels propagate at object granularity. A process
that reads one byte from a labeled file acquires the full label set. This
can produce over-tainting relative to byte-level
systems~\cite{kemerlis2012libdft}. Section~\ref{sec:evaluation} measures the
false-positive rate and validates that declassification and gate paths
prevent unnecessary harness intervention.
